\section{Ficha de Requisitos Legales para el Tratamiento de Datos}
\apendiceindice{F}{Ficha de Requisitos Legales para el Tratamiento de Datos}
La ficha reúne lo que el prototipo debe cumplir para tratar datos personales
de los estudiantes durante el piloto. Se llena con la revisión de la ley
estatal \parencite{edomex2017datos} y de la ley general
\parencite{lgpdppso2025} de protección de datos personales en posesión de
sujetos obligados, y se completa con la entrevista a Control Escolar
(Apéndice~B, pregunta 4). Al revisar cada ley se anota el artículo que
sustenta cada requisito.

\medskip
\noindent
{\small\linespread{1}\selectfont\renewcommand{\arraystretch}{1.4}
	\begin{tabular}{|p{3cm}|p{5.2cm}|p{2.6cm}|p{3.6cm}|}
		\hline
		\textbf{Requisito} & \textbf{Cómo se atiende en el prototipo} & \textbf{Artículo y ley} & \textbf{Acuerdo con Control Escolar} \\
		\hline
		Aviso de privacidad & Se entrega antes del alta del estudiante (Apéndice~G). & & \\[1.2em] \hline
		Consentimiento & Firma del estudiante antes del alta; quien no firme entra con la revisión manual. & & \\[1.2em] \hline
		Finalidad & Los datos se usan solo para validar el ingreso y evaluar el prototipo. & & \\[1.2em] \hline
		Proporcionalidad & Solo se registran matrícula, nombre, carrera, turno, fecha y hora del intento, dispositivo y resultado; no se usan datos biométricos. & & \\[1.2em] \hline
		Seguridad & Comunicación por HTTPS, códigos firmados con HMAC y acceso al panel con usuario y contraseña. & & \\[1.2em] \hline
		Conservación y supresión & Los registros se eliminan al terminar el periodo acordado. & & \\[1.2em] \hline
		Derechos de acceso, rectificación, cancelación y oposición & El estudiante puede pedir sus registros o su baja del piloto con el responsable designado. & & \\[1.2em] \hline
		Responsable & Área del plantel que resguarda los registros del piloto. & & \\[1.2em] \hline
	\end{tabular}}
